% 来源：Text2Mem，PDF第1页；Findings of ACL 2026，2105--2119
\papersection{Text2Mem}{统一的记忆操作语言与状态验证}{ngocnamd93-spec}

% 来源：原论文PDF第1--3页，图1；动机案例为教学概述
\begin{frame}{记忆操作的目标、动作与效果}
    \begin{block}{用户说：“暂时不要再提午餐。”}
        降低显著性、设置到期规则与删除记录，会产生不同状态。
    \end{block}
    \vspace{2mm}
    \begin{itemize}\setlength{\itemsep}{3mm}
        \item 检索解决“找回什么”；记忆控制还要明确动作、范围与时间。
        \item 输入是自然语言请求与上下文；输出是操作Schema或有序操作列表。
        \item Text2Mem提供操作规范与基准，主实验在SQL/SQLite参考后端进行。
    \end{itemize}
    \takeaway{明确目标与动作，检查状态变化}
    \note{约50秒。用户说暂时不要再提午餐，可能希望降低优先级、设置到期时间，或者删除记录，这些动作会产生不同状态。检索负责找回信息，记忆操作还要明确改变哪些记录、怎样改变。Text2Mem用结构化语言表达这些要求，并在SQL参考后端检查执行效果。}
\end{frame}

% 来源：原论文PDF第4页表1、第4--5页§3.2；附录A第11页
\begin{frame}{十二种操作与五个字段}
    \small
    \begin{tabularx}{\linewidth}{@{}p{20mm}Y@{}}
        \toprule
        阶段 & 操作 \\
        \midrule
        编码 ENC & Encode \\
        存储 STO & Update、Label、Promote、Demote、Merge、Split、Delete、Lock、Expire \\
        检索 RET & Retrieve、Summarize \\
        \bottomrule
    \end{tabularx}\par
    \vspace{3mm}
    \normalsize
    \key{stage / op / target / args / meta}\par
    \vspace{2mm}
    阶段、动作、目标、参数、执行上下文；各操作的必填条件不同。
    \par\vspace{3mm}target 支持 ids、filter、search 或 all。\par
    Demote调整优先级，Delete删除记录，Expire设置到期规则。\par
    \note{约70秒。十二种操作按编码、存储和检索三个阶段组织。Encode写入新记忆，存储阶段包含更新、合并、拆分、删除等动作，检索阶段包含Retrieve和Summarize。五个字段分别说明阶段、动作、目标、参数和上下文。其中Demote降低优先级，Delete删除记录，Expire设置到期规则。定位已有条目可以按ID、条件或搜索结果选择。}
\end{frame}

% 来源：依据原论文PDF第3页图2、第5页§3.3、第6页图3重绘；非运行记录
\begin{frame}{从操作生成到状态核验}
    \centering
    \begin{tikzpicture}[node distance=5mm and 8mm,>=Latex,
        box/.style={draw=accent,rounded corners=1mm,fill=accent!5,
        text width=34mm,minimum height=10mm,align=center,font=\small}]
        \node[box] (plan) {自然语言与上下文\\模型生成操作};
        \node[box,right=of plan] (val) {Validator\\结构、跨字段约束};
        \node[box,right=of val] (parser) {Parser\\类型化、参数归一化};
        \node[box,below=of parser] (adapter) {Adapter\\映射后端动作};
        \node[box,left=of adapter] (state) {状态变化与\\ExecutionResult};
        \node[box,left=of state] (check) {SQL断言\\检查预期效果};
        \draw[->,thick,accent] (plan)--(val);
        \draw[->,thick,accent] (val)--(parser);
        \draw[->,thick,accent] (parser)--(adapter);
        \draw[->,thick,accent] (adapter)--(state);
        \draw[->,thick,accent] (state)--(check);
    \end{tikzpicture}
    \vspace{3mm}\par
    \key{格式合法、动作可执行、效果符合意图，需要分别判断}
    \par\vspace{3mm}锁定、到期与字段约束由校验器检查；目标理解依赖模型。\par
    \note{约65秒。模型先生成操作，Validator检查结构和跨字段约束，Parser统一类型与参数，Adapter把操作映射到具体后端。ExecutionResult记录影响的条目和状态，最后用SQL断言核对预期效果。格式检查、执行检查和状态检查各有分工；适配新后端时，需要实现相应动作及约束。}
\end{frame}

% 来源：改写自原论文PDF第12页附录B.1.1；本页是未执行的教学案例
\begin{frame}{操作示例：提高进度笔记的权重}
    \small
    \begin{block}{请求与已明确的目标}
        “提高本周课设进度笔记的重要性，其他记录不变。”目标ID：17。
    \end{block}
    \begin{columns}[T,onlytextwidth]
        \column{.45\textwidth}
        \key{结构化操作}\par\vspace{2mm}
        \texttt{stage: STO}\par
        \texttt{op: Promote}\par
        \texttt{target: ids=[17]}\par
        \texttt{args: weight=0.9}\par

        \column{.51\textwidth}
        \key{执行后检查}\par\vspace{2mm}
        \begin{enumerate}\setlength{\itemsep}{2.5mm}
            \item ID 17的权重是否达到要求？
            \item ID 17的内容是否保持不变？
            \item 非目标条目是否保持不变？
        \end{enumerate}
        \vspace{2mm}
        若最初选错目标，合法操作也会改错记录。
    \end{columns}
    \note{约70秒。这个示例把目标设为编号17的课设进度笔记，要求提高权重，并保持其他记录不变。Promote操作把示例权重设为0.9。执行后有三项检查：目标权重达到要求，目标内容没有改变，其他记录也保持原样。如果目标在最初就选错了，后续仍可能执行成功，因此目标定位也要纳入检查。}
\end{frame}

% 来源：原论文PDF第6--8页§4.1--4.3；EMR公式在第7页
\begin{frame}{结构、执行与状态评价}
    \small
    \begin{tabularx}{\linewidth}{@{}p{22mm}Y@{}}
        \toprule
        指标 & 回答的问题 \\
        \midrule
        SMA & 操作与参考结构有多相似？使用编辑距离或余弦相似度。 \\
        ESR & 操作能否通过校验，并无运行错误地执行？ \\
        EMR & 全部期望SQL断言中，有多少得到满足？ \\
        \bottomrule
    \end{tabularx}\par
    \vspace{2mm}
    \[
        \mathrm{EMR}=\frac{\text{满足的状态断言数}}{\text{期望状态断言总数}}
    \]
    \par\vspace{3mm}评测覆盖五个模型，以及中英文、直接／间接指令、单操作／工作流。\par
    \note{约60秒。基准先建立记忆初态，再让模型生成操作并执行。SMA比较操作结构，ESR统计能够通过校验并正常执行的操作，EMR统计满足的状态断言。一个操作可能对应多条断言，因此ESR和EMR的分母不同。用这三项指标，可以分别观察表达、执行和状态效果中的问题。}
\end{frame}

% 来源：原论文PDF第8页表2 Avg列、表3；全部为作者结果
\begin{frame}{执行成功率与状态匹配结果}
    \small
    \begin{tabularx}{\linewidth}{@{}Yrr@{}}
        \toprule
        模型 & ESR（\%） & EMR（\%） \\
        \midrule
        GPT-4o & 94.84 & 25.32 \\
        Qwen2.5-72B-Instruct & 95.10 & 25.35 \\
        Claude-4-Sonnet & 95.94 & 31.32 \\
        Gemini-2.5-Pro & 94.79 & 28.97 \\
        GPT-5 & 96.62 & 28.26 \\
        \bottomrule
    \end{tabularx}\par
    \vspace{3mm}
    \normalsize\key{GPT-5工作流：ESR 100\%，EMR 11.90\%}\par
    \vspace{2mm}\small

    \par\vspace{3mm}该基准中，五个模型在 Split 操作上的 EMR 均为0。\par
    \note{约90秒。五个模型的ESR都接近百分之九十五，EMR则在百分之二十五到三十一左右。GPT-5的ESR最高，Claude的EMR最高。工作流中，GPT-5的操作全部可以执行，但只满足百分之十一点九的状态断言。Split也是困难项，五个模型在该基准中的EMR都为零。这些结果说明，多步执行和状态约束仍是主要问题。}
\end{frame}

% 来源：原论文PDF第9页Limitations、第15页附录F；课设联系为阅读分析
\begin{frame}{方法局限与课设启发}
    \begin{columns}[T,onlytextwidth]
        \column{.47\textwidth}
        \key{论文价值与局限}\par\vspace{2mm}
        \small
        \begin{itemize}\setlength{\itemsep}{2.5mm}
            \item 统一动作、范围和生命周期参数。
            \item 规划与状态效果分开评测。
            \item 缺少复杂语义保证；跨真实后端和长期任务证据有限。
        \end{itemize}
        \column{.48\textwidth}
        \key{与本组课设的联系}\par\vspace{2mm}
        \small
        \begin{itemize}\setlength{\itemsep}{2.5mm}
            \item StructMem研究事件记忆与问答，任务和指标不同。
            \item 在数据校验之外，增加操作后的状态检查。
            \item 借鉴显式输入、影响范围与非目标状态检查。
        \end{itemize}
    \end{columns}
    \takeaway{按目标记录与非目标记录分别核验操作效果}
    \note{约60秒。Text2Mem把动作、范围和生命周期参数写成明确的操作格式，再分别评价规划与状态效果。当前证据主要来自SQL参考环境，复杂语义、其他后端和长期任务仍需进一步研究。对我们的课设，可在数据校验之外增加操作后的状态检查，尤其关注目标范围和非目标记录。下面进入共同总结，比较六篇分别解决了哪些问题。}
\end{frame}
